\section{Related Work}
\label{sec:related}

\subsection{Rule-Based and Learned Intrusion Detection}
Signature-oriented network security systems emphasize transparency and deterministic protocol analysis. Snort operationalized lightweight signature matching~\citep{roesch1999snort}, while Bro/Zeek demonstrated stateful network analysis and policy-driven event processing~\citep{paxson1999bro}. These systems remain important because their behavior is auditable and known-pattern false positives can be controlled. Their limitation is coverage: semantic variants, authorization-sensitive requests, and business-logic abuse may not be reducible to a fixed payload pattern.

Machine-learning systems broaden the decision boundary beyond explicit signatures. DeepLog models execution logs as sequences for anomaly detection~\citep{du2017deeplog}; LUCID presents a lightweight deep-learning approach to DDoS classification~\citep{doriguzzi2020lucid}; and contextual encoders have been applied to network and intrusion-detection features~\citep{yang2025bertids}. Surveys of explainable intrusion detection note recurring trade-offs between broader pattern recognition, interpretability, and distributional stability~\citep{neupane2022xids}. The present study focuses on request-side HTTP evidence and a typed security decision rather than a single anomaly score.

\subsection{Local Language Models for Cybersecurity Analysis}
Language models are increasingly studied as aids for security triage, log analysis, vulnerability interpretation, and network-traffic reasoning. Their appeal is that heterogeneous evidence can be represented as text and mapped to structured explanations or decisions. Their weakness is that fluent output can overstate certainty, and models may respond differently to minor changes in prompt structure or data representation.

The closest recent empirical precedent is the evaluation of local LLMs for network traffic packet analysis by \citet{paukkeri2026local}. That study examined Llama 3.1 in depth and compared five local models on IoT packet data. It found that models recognized some language-based malicious indicators, especially SQL-related content, but were unreliable across attack types, benign traffic, hexadecimal input, and encrypted payloads. The present work shares the reliability-oriented perspective but differs in five ways. First, it studies request-level HTTP security decisions rather than individual packet interpretation. Second, it evaluates a task-adapted compact model against rules, conventional machine learning, untuned prompts, merged artifacts, and multiple serving paths. Third, it uses a frozen external challenge set that is isolated from training and tuning. Fourth, it treats insufficient context as a scored \texttt{unknown} class. Fifth, it applies paired confidence intervals and significance tests to same-sample predictions.

The contribution is therefore not another claim that an LLM can recognize obvious attack strings. It is a controlled estimate of how much task adaptation changes decisions, which failure modes remain, and whether those changes support a defensible security role.

\subsection{Parameter-Efficient Adaptation and Structured Decisions}
Transformers provide the underlying sequence-modeling architecture~\citep{vaswani2017attention}. The evaluated base model is Qwen3-4B-Instruct-2507~\citep{qwen3,qwen2507}. LoRA freezes pretrained weights and learns low-rank updates~\citep{hu2021lora}; QLoRA combines low-rank adapters with a frozen 4-bit quantized base model~\citep{dettmers2023qlora}. These techniques make local specialization feasible, but they do not guarantee calibrated uncertainty, safe false-positive behavior, or domain transfer.

This paper uses QLoRA as an established adaptation mechanism, not as a methodological novelty. The research question is empirical: whether adaptation improves attack-oriented HTTP classification under a sealed protocol, and whether the resulting error profile remains compatible with security decision support.

\subsection{Prompt Injection and Untrusted Request Content}
Prompt injection exploits the inability of instruction-following models to reliably separate trusted control text from untrusted content. PromptInject demonstrated goal hijacking and prompt leakage~\citep{perez2022promptinject}; indirect prompt injection extended the threat to retrieved or tool-supplied content~\citep{greshake2023indirect}; and HackAPrompt documented diverse bypass strategies~\citep{schulhoff2023hackaprompt}. In an HTTP classifier, the body, path, query string, and some headers are attacker-controlled data. The model must analyze those fields without obeying embedded instructions. The study therefore preserves raw outputs and evaluates message structure, schema compliance, and prompt-injection cases separately from verdict accuracy.

\subsection{Evaluation, Abstention, and Reproducibility}
Security evaluation is vulnerable to leakage, prevalence distortion, and selective metric reporting. We use stratified bootstrap confidence intervals~\citep{efron1993bootstrap}, paired McNemar tests~\citep{mcnemar1947note}, and Holm--Bonferroni correction~\citep{holm1979sequentially}. In addition to accuracy and F1, we report safe false-positive rate and unknown-class recall. The unknown class represents cases in which request-visible evidence is insufficient for a justified safe/attack decision; operationally, it functions as an abstention target.

A model can fail by choosing the wrong verdict, by following attacker-controlled instructions, by producing an invalid schema, or by expressing unjustified certainty. Reporting only aggregate accuracy collapses these failure channels. The present work scores all samples and reports each channel explicitly, while also testing whether artifact format and message serialization alter decisions.
